\documentclass[12pt]{exam}

\begin{document}

% =============================
% SUJET A
% =============================

\begin{center}
\Large BTS CIEL2 — QCM Langage C — Sujet A
\end{center}
\begin{center}
12/11/2025 \hfill Nom et prénom:\enspace \hrulefill
\end{center}

\textbf{Cochez la ou les bonnes réponses :}

\begin{questions}

\question Quelle instruction permet d’obtenir la taille du tableau \texttt{int tab[]} :
\begin{checkboxes}
 \choice \texttt{sizeof(tab)}
 \choice \texttt{len(tab)}
 \CorrectChoice \texttt{sizeof(tab) / sizeof(int)}
\end{checkboxes}

\question Quelles sont les caractéristiques du typage en C :
\begin{checkboxes}
 \CorrectChoice statique et faible
 \choice dynamique et faible
 \choice dynamique et strict
\end{checkboxes}

\question Quelle est la différence entre \texttt{calloc} et \texttt{malloc} :
\begin{checkboxes}
 \choice Elles sont identiques
 \CorrectChoice \texttt{calloc} initialise la mémoire à zéro
 \choice \texttt{malloc} initialise la mémoire avec des valeurs aléatoires
\end{checkboxes}

\question Lors de l’allocation dynamique, la mémoire est :
\begin{checkboxes}
 \CorrectChoice sur le tas
 \choice sur la pile
 \CorrectChoice gérée via un pointeur
\end{checkboxes}

\question Les mots clés \texttt{signed} et \texttt{unsigned} servent à :
\begin{checkboxes}
 \CorrectChoice définir la représentation d’un entier (positif/négatif)
 \choice choisir si le compilateur doit signer le code
 \choice typer une variable Unicode
\end{checkboxes}

\question Résultat de \texttt{11 / 2} :
\begin{checkboxes}
 \choice 5.5
 \CorrectChoice 5
 \choice comportement indéfini
\end{checkboxes}

\question Que signifie \texttt{int *p = NULL;} :
\begin{checkboxes}
 \choice Pointe vers la case mémoire 0
 \CorrectChoice Le pointeur ne pointe vers aucune adresse valide
 \choice Le pointeur contient une adresse non initialisée
\end{checkboxes}

\pagebreak{}
\question Quelle syntaxe est invalide pour une boucle \texttt{for} :
\begin{checkboxes}
 \CorrectChoice \texttt{for(i in list)\{ \}}
 \choice \texttt{for(int i=0; i<3; i++);}
 \choice \texttt{for(;;)\{ \}}
\end{checkboxes}

\question Quelle valeur est fausse au sens booléen :
\begin{checkboxes}
 \choice 1
 \CorrectChoice 0
 \choice "" (chaîne vide)
\end{checkboxes}

\question À quoi sert une structure :
\begin{checkboxes}
 \CorrectChoice Structurer les données en mémoire
 \CorrectChoice Type composite regroupant plusieurs champs
 \choice Construire ses propres structures de contrôles
\end{checkboxes}

\question Si une zone allouée n’est jamais libérée :
\begin{checkboxes}
 \choice Libération automatique à la fin du programme
 \CorrectChoice Fuite mémoire
 \choice Allocation réutilisée par le système
\end{checkboxes}

\question Fonctions pour l’allocation dynamique :
\begin{checkboxes}
 \choice mem\_alloc, mem\_free, mem\_init
 \CorrectChoice malloc, free, realloc
 \choice new, delete, sizeof
\end{checkboxes}

\question Lequel de ces types utilise le moins de mémoire :
\begin{checkboxes}
 \choice signed short
 \CorrectChoice char
 \choice unsigned float
\end{checkboxes}

\question Quelle est la conséquence d’un \texttt{free()} multiple :
\begin{checkboxes}
 \choice Ignoré par le système
 \CorrectChoice Comportement indéfini
 \choice Libère seulement la dernière occurrence
\end{checkboxes}

\question Où sont stockées les variables locales :
\begin{checkboxes}
 \choice En mémoire globale
 \CorrectChoice Sur la stack
 \choice Sur la heap
\end{checkboxes}

\end{questions}

\newpage
\setcounter{page}{1}

% =============================
% SUJET B
% =============================

\begin{center}
\Large BTS CIEL2 — QCM Langage C — Sujet B
\end{center}
\begin{center}
12/11/2025 \hfill Nom et prénom:\enspace \hrulefill
\end{center}

\textbf{Cochez la ou les bonnes réponses :}

\begin{questions}

\question Quelle définition décrit un tableau :
\begin{checkboxes}
 \choice Ensemble d’éléments de taille variable
 \CorrectChoice Zone mémoire contiguë d’éléments du même type
 \CorrectChoice Adresse mémoire constante
\end{checkboxes}

\question Quelle est la différence entre la stack et la heap :
\begin{checkboxes}
 \CorrectChoice Stack = locales, Heap = dynamiques
 \choice Stack = manuel, Heap = auto
 \CorrectChoice Stack = auto, Heap = manuel
\end{checkboxes}

\question Pourquoi libérer la mémoire :
\begin{checkboxes}
 \choice Ce n'est plus nécessaire avec les systèmes récents
 \CorrectChoice Pour éviter une fuite mémoire
 \choice Cela évite au Garbage Collector de se déclencher
\end{checkboxes}

\question Quelles sont les valeurs considérées comme fausses :
\begin{checkboxes}
 \CorrectChoice 0
 \choice "" (chaîne vide)
 \choice 1
\end{checkboxes}

\question Qu’est-ce qu’un pointeur \texttt{void *} :
\begin{checkboxes}
 \choice Pointeur vers une fonction
 \CorrectChoice Pointeur générique sans type
 \choice Pointeur vers une adresse constante
\end{checkboxes}

\question Quelles sont les caractéristiques du typage en C :
\begin{checkboxes}
 \choice dynamique et strict
 \CorrectChoice statique et faible
 \choice dynamique et faible
\end{checkboxes}

\question Code invalide :
\begin{checkboxes}
 \CorrectChoice \texttt{for(i in list)\{ \}}
 \choice \texttt{for(int i=0; i<3; i++);}
 \choice \texttt{for(;;)\{ \}}
\end{checkboxes}

\pagebreak{}
\question Que fait \texttt{sizeof(char)} :
\begin{checkboxes}
 \choice Retourne 8
 \CorrectChoice Retourne 1
 \choice Taille dépend du compilateur
\end{checkboxes}

\question Erreur courante avec les chaînes :
\begin{checkboxes}
 \CorrectChoice Buffer Overflow
 \choice Dirty Stack
 \choice Character Spoofing
\end{checkboxes}

\question Quel est le comportement attendu avec ce code :
\begin{verbatim}
int *cool_function() {
    int a = 0;
    return &a;
}
\end{verbatim}
\begin{checkboxes}
 \choice Stack overflow
 \CorrectChoice Segmentation Fault
\end{checkboxes}

\question Quand \texttt{malloc} retourne \texttt{NULL} :
\begin{checkboxes}
 \CorrectChoice Quand la mémoire n’a pas pu être allouée
 \choice Jamais
 \choice Lorsque le pointeur initial vaut NULL
\end{checkboxes}

\question Quelle est la conséquence d’un oubli de \texttt{free()} :
\begin{checkboxes}
 \choice Comportement indéfini
 \CorrectChoice Fuite mémoire
 \choice Crash du système
\end{checkboxes}

\question Que fait \texttt{strcpy(dest, src);} :
\begin{checkboxes}
 \choice Affecte un pointeur vers \texttt{src}
 \CorrectChoice Copie la chaîne \texttt{src} dans \texttt{dest}
\end{checkboxes}

\question À quoi sert le mot-clé \texttt{struct} :
\begin{checkboxes}
 \choice Créer un tableau dynamique
 \CorrectChoice Définir un type de données composite
 \choice Créer sa propre structure de contrôle
\end{checkboxes}

\question Différence entre \texttt{malloc} et \texttt{realloc} :
\begin{checkboxes}
 \choice \texttt{malloc} met la mémoire à zéro
 \CorrectChoice \texttt{realloc} utilise un bloc existant
 \choice realloc est obsolète
\end{checkboxes}

\end{questions}

\end{document}
